% 第 8 题燃油效率三线图：按固定源内容流矢量层精确值复刻（1 pt → 0.03515 cm，与全库同制）。
% 原点取坐标系原点（0 km/h, 0 km/L）= 源页 (508.4940, 559.1030) pt，y 向上；
% 刻度换算：x 0.992135 pt/(km/h)、y 7.93708 pt/(km/L)（由刻度线 39.6854 pt 间距反推）。
% 源实测：x 轴至 663.6490（箭头顶点 667.2350）、y 轴至 423.7910（箭头顶点 420.2043）；
% 刻度线 40/80（x，向上 3.969）与 5/10/15（y，向右 3.969）——照源坐标区间（559.1030→555.1340／508.4940→512.4630）；
% 修复轮 1（2026-10-10）：首版误采信 figure-inventory §4.2 表头散文「向下/向左」，与同表坐标区间相反；已镜像归正。
% 三系列：甲＝虚线[2.989 2.989]＋实心圆点×10；乙＝点线[0.398 1.993]＋实心三角×15；丙＝实线＋实心方块（2.1×2.1）×14。
% 末点标记兼作图例键，甲车/乙车/丙车文字紧排其右（x 起点一律源 643.15）；74 km/h 越格点照源。
\begin{tikzpicture}[x=0.03515cm,y=0.03515cm,>=stealth,every node/.style={font=\small,inner sep=0pt},
  gl/.style={line width=0.398bp,line cap=round,line join=round}]
% ===== 坐标轴与刻度 =====
\draw[gl,->] (0,0)--(158.7410,0);        % x 轴（线段至 155.1550，箭头顶点照源 158.7410）
\draw[gl,->] (0,0)--(0,138.8987);        % y 轴（线段至 135.3120，箭头顶点照源 138.8987）
\draw[gl] (39.6854,0)--(39.6854,3.9690);
\draw[gl] (79.3703,0)--(79.3703,3.9690);
\draw[gl] (0,39.6854)--(3.9690,39.6854);
\draw[gl] (0,79.3707)--(3.9690,79.3707);
\draw[gl] (0,119.0561)--(3.9690,119.0561);
\node at (-6.3190,-7.8420) {0};
\node at (39.6860,-7.8420) {40};
\node at (79.3760,-7.8420) {80};
\node at (-6.3190,38.9280) {5};
\node at (-8.9440,78.6130) {10};
\node at (-8.9440,118.2980) {15};
% ===== 甲车（虚线＋实心圆点）=====
\draw[gl,dash pattern=on 2.989pt off 2.989pt]
  (9.9213,39.6853)--(19.8423,71.4333)--(33.7323,99.2133)--(39.6853,119.0563)--
  (59.5283,103.1823)--(79.3703,79.3703)--(89.2923,95.2443)--(99.2133,91.2763)--
  (109.1343,103.1823)--(130.9613,79.3703);
\foreach \x/\y in {9.9213/39.6853, 19.8423/71.4333, 33.7323/99.2133, 39.6853/119.0563,
  59.5283/103.1823, 79.3703/79.3703, 89.2923/95.2443, 99.2133/91.2763,
  109.1343/103.1823, 130.9613/79.3703} \fill (\x,\y) circle[radius=1.575];
% ===== 乙车（点线＋实心三角）=====
\draw[gl,dash pattern=on 0.398pt off 1.993pt]
  (7.9369,11.9057)--(15.8739,15.8737)--(23.8109,23.8117)--(31.7479,27.7797)--
  (39.6849,35.7167)--(47.6219,39.6857)--(55.5589,43.6537)--(63.4969,45.6377)--
  (73.4179,47.6227)--(79.3709,48.8127)--(87.3079,55.5597)--(95.2449,61.9087)--
  (109.1349,59.5277)--(123.0249,51.5907)--(130.9619,49.6067);
\foreach \x/\y in {7.9369/11.9057, 15.8739/15.8737, 23.8109/23.8117, 31.7479/27.7797,
  39.6849/35.7167, 47.6219/39.6857, 55.5589/43.6537, 63.4969/45.6377,
  73.4179/47.6227, 79.3709/48.8127, 87.3079/55.5597, 95.2449/61.9087,
  109.1349/59.5277, 123.0249/51.5907, 130.9619/49.6067}
  \fill (\x,\y+2.2)--(\x-1.9,\y-1.1)--(\x+1.9,\y-1.1)--cycle;
% ===== 丙车（实线＋实心方块）=====
\draw[gl]
  (7.9369,19.8426)--(15.8739,23.8116)--(23.8109,39.6856)--(31.7479,55.5596)--
  (39.6849,49.6066)--(47.6219,47.6226)--(55.5589,55.5596)--(63.4969,59.5276)--
  (73.4179,54.3686)--(79.3709,52.7816)--(87.3079,50.7976)--(109.1349,47.6226)--
  (123.0249,35.7166)--(130.9619,31.7486);
\foreach \x/\y in {7.9369/19.8426, 15.8739/23.8116, 23.8109/39.6856, 31.7479/55.5596,
  39.6849/49.6066, 47.6219/47.6226, 55.5589/55.5596, 63.4969/59.5276,
  73.4179/54.3686, 79.3709/52.7816, 87.3079/50.7976, 109.1349/47.6226,
  123.0249/35.7166, 130.9619/31.7486}
  \fill (\x-1.05,\y-1.05) rectangle (\x+1.05,\y+1.05);
% ===== 轴题与图例文字（照源字形包围盒中心定位；括号按全库先例用全角）=====
\node at (24.6960,139.9130) {燃油效率};
\node at (66.3310,139.3430) {（km/L）};
\node at (140.2160,-7.9770) {速度};
\node at (170.9910,-8.5370) {（km/h）};
\node at (145.1560,79.9530) {甲车};
\node at (145.1560,50.1830) {乙车};
\node at (145.1560,32.3230) {丙车};
\end{tikzpicture}
